\documentclass[11pt,a4paper]{article}
\usepackage[T1]{fontenc}
\usepackage{amsmath,amssymb,geometry}
\geometry{margin=27mm}
\title{EPPA--III: excluding mixed partial profiles by saturation}
\author{Independent validation supplement; manuscript frozen}
\date{11 October 2026}
\begin{document}
\maketitle

\paragraph{Statement.}
Let \(G=sK_t\), where \(s,t\ge2\), be an induced subgraph
of a finite EPPA witness \(H\) with \(|H|=2st\).
Let \(w\in H\setminus G\), let \(P_i=N_H(w)\cap C_i\),
and call coordinate \(i\) \emph{partial} if
\(0<|P_i|<t\). If \(w\) has two or more partial
coordinates, necessarily \(s=t=2\), and both partial
coordinates have size \(1\).

\paragraph{Proof (using the already verified Claim 3.9 bound).}
Every automorphism of \(G\) extends to an automorphism of \(H\),
so the full \(W=S_t\wr S_s\) orbit of \(P=(P_i)_{i=1}^s\)
has size at most \(st\).
The Lean-verified Claim~3.9 estimate says that
each partial coordinate has size \(1\) or \(t-1\).
If there are \(r\ge2\) coordinates of the same partial
size \(k\), then the already Lean-verified lower bound gives
\[
 \binom sr\binom tk^r=\binom sr t^r
 \ \le\ |W\cdot P|\ \le\ st.
 \tag{1}
\]
The arithmetic theorem
\texttt{multiple\_partial\_orbit\_arithmetic}
forces \(s=t=r=2\).
For \(t=2\) the only partial size is \(k=1\),
so this handles all multiple-partial profiles.

It remains to suppose \(t\ge3\) and that two
partial coordinates have \emph{different} sizes.
By the preceding paragraph, there must be exactly
one singleton coordinate and exactly one
co-singleton coordinate. Let \(i\) be the singleton
coordinate.
For any index \(j\) and singleton \(A\subseteq C_j\),
an outer permutation taking \(i\) to \(j\), followed
by an inner permutation, produces a profile \(Q\)
in \(W\cdot P\) whose unique singleton coordinate is
\(j\) and whose singleton is \(A\).
Thus the map
\[
 \Phi:W\cdot P\longrightarrow
 \{(j,A):j\in[s],\ A\in\binom{C_j}{1}\},
 \qquad
 Q\longmapsto(\hbox{its singleton index and subset})
\]
is surjective. The target has \(st\) elements,
whereas the source has at most \(st\), so \(\Phi\)
must be injective.

However, choose a nontrivial permutation of the
vertices in the co-singleton coordinate that
moves its \(t-1\)-subset, acting identically on all
other coordinates. This gives \(Q'\ne P\) in the
same wreath orbit with \(\Phi(Q')=\Phi(P)\),
a contradiction. Thus there are no mixed-size
multiple-partial profiles. \(\square\)

\paragraph{Consequence for Lemma 3.7.}
This replaces the need for the more general
multinomial estimate
\(\binom{s}{p}t^p\le st\) in the
direct classification: the existing single-size
wreath-orbit bound, an elementary saturation
argument, and independent inner permutations suffice.
With at most one partial coordinate, the orbit
of its profile has exactly
\(s\binom tk\binom{s-1}{r}\) elements, where \(r\)
is the number of other full coordinates.
The same \(st\) bound forces \(r=0\) or \(s-1\);
the orbit exhausts \(H\setminus G\).
The latter is then \(sK_t\) by the edge-count
identity under vertex-transitivity.
When \(s=t=2\) and there are two partial
coordinates, the four profiles yield precisely
the Wagner exception. Profiles without any partial
coordinates are handled by the true-twin quotient
argument in \texttt{direct-clique-cover-proof.tex}.

\paragraph{Formal status and scope.}
The full \(W\)-orbit lower bound for a \emph{fixed}
partial size, its transfer into the actual EPPA orbit,
the extreme-degree restriction, and the arithmetic
in (1) are Lean-verified (merged PRs~\#25, \#29, \#31).
The injectivity of the saturated choice projection is
the goal of PR~\#32; the final step selecting a
nontrivial permutation on a second partial coordinate
is a paper-level argument, not yet Lean-checked.
The original manuscript and its existing Wagner
patch are untouched.

\end{document}
